\section{引言}

\subsection{背景与动机}
可再生能源并入电网推动了虚拟电厂（VPPs）的发展，以聚合和管理分布式能源资源。对 VPP 动态——特别是电力负荷、光伏（PV）发电和风力发电——的准确预测对于电网稳定性、经济调度和运行风险管理至关重要 \cite{Ma2025,Cao2023,Qiu2024,Moreno2020,Wu2024}。与单变量预测不同，协作式 VPP 管理需要跨这些相互关联的领域进行同步、协调的预测，且必须严格遵守每种能源形式的物理约束。例如，光伏发电本质上受太阳辐射约束，在夜间无法发电，而风力发电则依赖于局部风速驱动的非线性风机功率曲线。

其早期预测方法多依赖于统计外推，而随着分布式能源的增加，Transformer 架构凭借其捕获长程时间动态的能力，已成为该领域的主流范式 \cite{Zhou2021informer,Wu2021autoformer,Nie2023patchtst,Zeng2023}。然而，可靠的调度决策不仅取决于最小化原始统计预测误差，还取决于预测模型是否从根本上遵守支配电力系统的物理定律。

\subsection{现有数据驱动方法的局限性}
尽管深度学习在时间序列预测中取得了进展，但其在 VPP 调度中仍面临三大核心挑战 \cite{Vle2025,Arabzadeh2025}：

\textbf{1) 纯数据驱动的物理越界：} 虽然神经网络在最小化统计误差方面表现出色，但其作为“黑盒”运行，忽视了物理原理。这导致模型经常产生物理上不可能的预测，如夜间非零太阳能或负值风电，严重损害了运行安全性和操作员信任。

\textbf{2) 物理信息增强（PINN）在随机环境下的僵化：} 传统的 PINN \cite{Raissi2019} 多将物理方程作为损失函数中的软约束。然而，VPP 运行在高度随机的环境中，固定参数的软约束难以动态适应多变的气象分布，且容易受到数据驱动项与物理惩罚项之间的梯度冲突影响，导致收敛困难 \cite{Li2026,Asinyo2025}。

\textbf{3) 结构性追溯缺失：} 现有模型缺乏因果可解释性。依赖事后归因或注意力可视化无法提供确定性的物理因果归因 \cite{Lovo2025}。电网操作员需要模型显式展示可测量环境变量（如辐照度）如何支配设备激活边界，而目前在动态关联网络激活与物理状态转换方面仍存在技术缺口 \cite{Song2025,Chen2024}。

\subsection{所提方法与主要贡献}
为此，本文提出了 \textbf{PhysFormer}，一种物理引导的因果 Transformer，通过将可微物理先验内生化，将神经网络锚定在内在的结构合规性上。其主要贡献如下：

1) **受限物理动作残差（BPAR）架构**：通过引入 Softplus 歧义消除与物理激活掩码，PhysFormer 确保预测轨迹在数学上驻留在绝对安全的流形上。该机制在不使用非连续性事后剪裁的情况下，将边界违规率（BVR）降至绝对 0.00\%，且能将边界处的结构性修正控制在物理可忽略的量级。

2) **正交物理引导因果耦合（PGCC）机制**：通过显式正则化对潜层门控进行物理对齐，实现了皮尔逊相关系数 $r=0.8306$ 的确定性物理关联。消融实验证实，这种透明性的获得与精度目标是正交的（即不牺牲 MSE 精度），同时使模型对 50\% 幅度的外源传感器噪声具备极强的韧性。


